\section{Introduction and main results}
\label{intro:section}

The random friend tree grows by selecting a uniformly random existing
vertex, selecting one of its neighbours uniformly, and attaching a new
leaf to that neighbour. Krapivsky and Redner introduced the model in
their study of emergent network modularity~\cite{KR}. Its number $C_n$
of non-leaf vertices grows on a nontrivial power scale. We study the
exponent $\mu$ characterized by
\[
 \mathbb E C_n=n^{\mu+o(1)}.
\]
The existence of this exponent and its identification with the
almost-sure logarithmic exponent are inputs from Uzay's
computer-assisted work~\cite{Uzay}. That work gives
$0.5650\le\mu\le0.5665$, excluding $17/30$.
The broader probabilistic structure of random friend trees is developed
in~\cite{AB}.

Our main questions concern exact characterization and effective
approximation. An infinite operator equation can identify a real number
without supplying a terminating algorithm for its approximation.
Likewise, numerical evidence of agreement between two initial
environments does not imply that their mean exponents agree uniformly
over all environments. We address both issues by controlling the
effect of a single boundary leaf on an arbitrary old component.

We embed the tree process in continuous time by assigning a rate-one
source clock to each real vertex. A \emph{core} is a vertex of full
degree at least two. In a retained component, the root has one inert
parent stub: it contributes one to the root's degree, supplies no
incoming leaves, and produces a null transition when selected. The
\emph{inert newborn} consists of a root core, one real leaf, and this
stub. Write $F_\infty(t)$ for its expected number of cores.

For a newborn core in a finite environment $x$, let $F_x(t)$ be the
expected number of its future core descendants, including itself, and
put
\[
 G(t)=\sup_xF_x(t),\qquad H(t)=1+\int_0^tG(u)\,du.
\]
The admissible class and its harmless finite inert-boundary completion
are specified in Section~\ref{model:section}. In the inert newborn
process, let $\sigma_v$ be the birth time of core $v$, and give the
initial core genealogical depth zero. Define
\[
 A_m^\infty(s)=\mathbb E_\infty
   \sum_{\operatorname{depth}(v)=m}e^{-s\sigma_v},\qquad s>0.
\]

\begin{theorem}[Exact characterization and effective computation]
\label{main:computation}
The limits
\[
 r_\infty(s)=\lim_{m\to\infty}A_m^\infty(s)^{1/m}
 \quad(s>0)
\]
exist and define a positive, continuous, strictly decreasing function.
Its unique crossing of one is the original friend-tree exponent:
\[
 r_\infty(\mu)=1.
\]
Moreover,
\[
 \lim_{t\to\infty}\frac{\log G(t)}t
 =\mu
 =\inf_{T\in\mathbb Q_{>0}}\frac{\log H(T)}T.
\]
Every fixed coefficient $A_m^\infty(s)$ at rational $s>0$, and every
$H(T)$ at rational $T>0$, is effectively computable. There is an
algorithm which, on input a positive rational $\varepsilon$, halts
with rational $a<b$ satisfying
\[
 a<\mu<b,\qquad b-a<\varepsilon.
\]
\end{theorem}

The proof combines Theorems~\ref{gc:root}, \ref{calc:calculator},
\ref{erc:equality}, and \ref{ecc:computable}. The lower certificate is
particularly explicit:
\[
 s<\mu\quad\Longleftrightarrow\quad
 A_m^\infty(s)>12(s+3)\quad\hbox{for some finite }m.
\]
Finite positive rational-history sums eventually witness every strict
lower inequality. The submultiplicative time-block formula supplies
the complementary complete search for strict upper inequalities.
The resulting algorithm has a proof of termination and a prescribed
error, but no useful complexity bound is established here.

For a retained state $y$, let $P_t^\infty$ be the inert semigroup and
write
\[
 V_s(y)=\int_0^\infty e^{-st}P_t^\infty C(y)\,dt,
 \qquad
 W(s)=\sup_x\int_0^\infty e^{-st}F_x(t)\,dt.
\]
If $y^+$ is obtained by adding a leaf at the retained root, the main
uniform estimate is as follows.

\begin{theorem}[Boundary gradient and robustness]
\label{main:boundary}
For every $s>\mu$, the preceding quantities are finite, and
\[
 |V_s(y^+)-V_s(y)|\le\frac{40}{s}W(s)
\]
uniformly over all finite retained states $y$. A finite time-block
certificate computes an explicit upper bound for the right side.
Fix $Q<\infty$ and permit an arbitrary nonanticipating root-leaf input
with predictable rate $0\le q_t\le Q$. Starting from the inert newborn,
let $F_Q^-(t)$ and $F_Q^+(t)$ be the infimum and supremum of the expected
core counts over these policies. Then
\[
 \lim_{t\to\infty}\frac{\log F_Q^-(t)}t
 =\lim_{t\to\infty}\frac{\log F_Q^+(t)}t=\mu.
\]
The conclusion includes policies selected separately for each horizon,
provided that the same bound $Q$ applies.
\end{theorem}

The finite-horizon gradient is proved before uniform resolvent
finiteness is known. Its proof uses a disagreement exploration with
discounted reproduction mass at most $19/20$ at $s=11/20$. This avoids
assuming the conclusion in the comparison. Theorems~\ref{erc:gradient}
and~\ref{bcl:matching} give the precise statements and proofs.

We also refine the existing rational motif certificate.

\begin{theorem}[A certified numerical refinement]
\label{main:bound}
The original exponent satisfies
\[
 \frac{113}{200}\le\mu\le
 \frac{2832242711}{5000000000}=0.5664485422.
\]
The lower endpoint is imported from~\cite{Uzay}. The upper endpoint
follows from an exact rational modification of its upper motif
potential, with a positive reserve in every vertex row.
\end{theorem}

This is a refinement of an existing computer-assisted certificate.
It reduces the interval width by approximately $3.43\%$; it does not
give high-precision digits of $\mu$. Section~\ref{numerical:section}
states the all-degree extension and the exact verification procedure.

\paragraph{Scope of the results.}
This paper proves mean exponential growth statements, exact
characterizations, and computability, using the stated exponent
existence input. It does not prove a sharp asymptotic
$\mathbb E C_n\sim c n^\mu$, a positive almost-sure amplitude, or
irrationality or transcendence of $\mu$. A separate companion draft
studies independent branching representations and irrationality for
modified local redirection models~\cite{Local}; those exponents are
distinct from the original exponent studied here.

\paragraph{Organization.}
Sections~\ref{model:section}--\ref{calc:section} give the exact kernel,
inert comparison, pressure equation, and finite coefficient calculator.
The uniform mean and locality estimates are followed by finite-window
supersolutions, controlled moments, and the bounded-parent likelihood
argument. Section~\ref{erc:section} proves the boundary gradient;
Section~\ref{erc:closure-section} identifies the uniform exponent.
The final sections derive computability, construct the resolvent,
prove boundary robustness, and give the numerical certificate.
